\section*{Ethical Considerations}

Inferring habits creates privacy risks beyond storing explicit user statements.
The system may derive routines, preferences, schedules, or other sensitive regularities that a user did not knowingly designate as memory.
Such profiles could reveal sensitive behavior if exposed, transferred across contexts, or accessed by an unauthorized party.
A deployment should therefore require informed opt-in and provide a visible habit inspector, correction, per-habit deletion, retention limits, and a global disable control.

Proactive interventions can also be unwanted, manipulative, mistimed, or unsafe.
Phase checks and per-habit cooldown constrain intervention frequency, while high-impact settings additionally require explicit confirmation and application-specific policy checks.
Assessing user trust, perceived helpfulness, and transfer to live use requires prospective evaluation beyond the offline labels.

Benchmark-derived prompts were sent to external model APIs during the experiments.
Released artifacts exclude credentials and personal paths and document transmitted fields and provider retention assumptions.
Audit traces can contain behavioral evidence and require the same access, retention, and deletion protections as the Habit Store.

